%% ============================================================
%% shared/services/compute/auto_scaling_elb.tex
%% EC2 Auto Scaling + Elastic Load Balancing — SAA + SAP
%% ============================================================

\servicesection{EC2 Auto Scaling}{\bothtag}{Automatically adjust EC2 fleet size based on demand}

\begin{keyfacts}
\begin{itemize}
  \item \textbf{Auto Scaling Group (ASG)}: logical group of EC2 instances; managed as a fleet.
  \item Defines: \textbf{Minimum}, \textbf{Desired}, and \textbf{Maximum} capacity.
  \item \textbf{Launch Template} (preferred) or Launch Configuration: specifies AMI, instance
        type, key pair, SG, IAM role, user data for new instances.
  \item Integrates with ELB: new instances are automatically registered with the target group;
        unhealthy instances are replaced.
  \item \textbf{Health checks}: EC2 status check (default) or ELB health check
        (more precise --- app-level).
  \item \textbf{AZ balance}: ASG distributes instances evenly across specified AZs.
\end{itemize}
\end{keyfacts}

\subsection{Scaling Policies}

\begin{comparebox}[title={Auto Scaling policy types}]
\begin{tabular}{@{}p{3.5cm}p{5cm}p{3.5cm}@{}}
\toprule
\textbf{Policy} & \textbf{How it works} & \textbf{Use case} \\
\midrule
\textbf{Target Tracking}  & Maintains a target value for a metric
    (e.g.\ CPU = 50\%) by adding/removing instances & Simplest; most common \\[4pt]
\textbf{Step Scaling}     & Triggers at alarm thresholds; scales by steps
    based on breach magnitude; does not wait for cooldown & Custom multi-step response \\[4pt]
\textbf{Simple Scaling}   & Single action per alarm; waits for cooldown
    before next action; legacy & Basic setups \\[4pt]
\textbf{Scheduled Scaling} & Pre-define capacity changes for known patterns
    (e.g.\ scale up at 8 AM Mon--Fri) & Predictable load \\[4pt]
\textbf{Predictive Scaling} & ML analyses historical patterns; provisions
    ahead of expected demand & Recurring traffic spikes \\
\bottomrule
\end{tabular}
\end{comparebox}

\begin{examtip}
\textbf{Target Tracking} is the default recommendation for most scenarios.
\textbf{Scheduled Scaling} is best when you know in advance (e.g.\ Black Friday launch,
business-hours pattern). \textbf{Predictive} + Target Tracking together is the SAP-level
answer for ``automatically handle cyclical load without manual intervention''.
\end{examtip}

\subsection{Lifecycle Hooks}

Pause instance launch or termination to perform custom actions:
\begin{itemize}
  \item \textbf{Launch hook}: wait for instance to complete bootstrapping before serving traffic
        (e.g.\ pull config from S3, run software install). Timeout: 1 hour.
  \item \textbf{Termination hook}: drain connections, push logs to S3, deregister from
        discovery service before instance is terminated.
\end{itemize}

\subsection{Multi-Tier Architecture Diagram}

\begin{center}
\begin{tikzpicture}[
  tier/.style={draw=awsdark, fill=awsdark!6, rounded corners=5pt,
               minimum width=3.8cm, minimum height=1.1cm,
               font=\small\bfseries, align=center},
  asg/.style={draw=saacolor, fill=saacolor!8, rounded corners=3pt,
              minimum width=2.6cm, minimum height=0.9cm,
              font=\tiny, align=center},
  db/.style={draw=sapcolor, fill=sapcolor!8, rounded corners=3pt,
             minimum width=2.6cm, minimum height=0.9cm,
             font=\tiny, align=center},
  arr/.style={->, thick, awsdark}
]

\node[tier] (dns)  at (0, 5.5)  {Route~53 DNS};
\node[tier] (cf)   at (0, 4.2)  {CloudFront};
\node[tier] (alb)  at (0, 2.9)  {ALB (Multi-AZ)};

\node[asg]  (web1) at (-2.0, 1.5) {Web / App\\ASG};
\node[asg]  (web2) at ( 2.0, 1.5) {Web / App\\ASG};
\node[font=\tiny, color=gray] at (0, 1.5) {AZ-a \quad AZ-b};

\node[db]   (rds)  at (-1.8, 0) {RDS Primary\\(AZ-a)};
\node[db]   (rds2) at ( 1.8, 0) {RDS Standby\\(AZ-b)};

\draw[arr] (dns) -- (cf);
\draw[arr] (cf)  -- (alb);
\draw[arr] (alb) -- (web1);
\draw[arr] (alb) -- (web2);
\draw[arr] (web1) -- (rds);
\draw[arr] (web2) -- (rds);
\draw[dashed, thick, saacolor] (rds) -- node[below, font=\tiny]{sync replication} (rds2);

\end{tikzpicture}
\end{center}

%% ────────────────────────────────────────────────────────────────

\servicesection{Elastic Load Balancing (ELB)}{\bothtag}{Distribute traffic across targets in one or more AZs}

\begin{keyfacts}
\begin{itemize}
  \item All ELB types are \textbf{highly available} by default (across AZs in the region).
  \item \textbf{Target types}: EC2 instances, IP addresses, Lambda functions, other ALBs (ALB only).
  \item \textbf{Health checks}: ELB removes unhealthy targets from rotation.
  \item \textbf{Cross-zone load balancing}: distributes traffic evenly across all registered
        targets regardless of AZ. Default on for ALB (cost per AZ for NLB/GWLB).
  \item \textbf{Sticky sessions (session affinity)}: route requests from same client to same
        target (ALB: cookie-based; NLB: source IP).
\end{itemize}
\end{keyfacts}

\subsection{Load Balancer Types Comparison}

\begin{comparebox}[title={ELB type selection}]
\begin{tabular}{@{}p{1.8cm}p{2.5cm}p{3cm}p{4cm}@{}}
\toprule
\textbf{Type} & \textbf{Layer} & \textbf{Protocols} & \textbf{Best for} \\
\midrule
\textbf{ALB}  & L7 (application) & HTTP, HTTPS, WebSocket &
  Host/path-based routing, microservices, containers, Lambda, gRPC \\[4pt]
\textbf{NLB}  & L4 (transport) & TCP, UDP, TLS &
  High-perf (millions req/s), gaming, IoT, static IP needed, PrivateLink \\[4pt]
\textbf{GWLB} & L3 (gateway) & GENEVE (IP packets) &
  Inline inspection via 3rd-party virtual appliances (NGFW, IDS/IPS) \\[4pt]
\textbf{CLB}  & L4/L7 & HTTP, HTTPS, TCP &
  Legacy; avoid for new workloads \\
\bottomrule
\end{tabular}
\end{comparebox}

\begin{examtip}
\begin{itemize}
  \item \textbf{``Content-based routing''} (host header, URL path, HTTP method) $\Rightarrow$ \textbf{ALB}.
  \item \textbf{``Static IP'' or ``millions of requests per second'' or UDP} $\Rightarrow$ \textbf{NLB}.
  \item \textbf{``Inspect all traffic through a firewall appliance''} $\Rightarrow$ \textbf{GWLB}.
  \item NLB supports \textbf{PrivateLink} (expose service to other VPCs/accounts) --- ALB does not.
  \item ALB supports \textbf{Lambda} as a target --- NLB does not.
\end{itemize}
\end{examtip}

\subsection{ALB Advanced Features}

\begin{itemize}
  \item \textbf{Listener rules}: route by host header (\texttt{api.example.com}), URL path
        (\texttt{/images/*}), query string, HTTP method, IP source. Priority-based evaluation.
  \item \textbf{Target groups}: pool of targets; each rule forwards to a target group.
        Can mix instance, IP, and Lambda targets.
  \item \textbf{Blue/green with ALB}: use weighted target groups to shift traffic gradually
        between old and new deployments.
  \item \textbf{Connection draining (deregistration delay)}: default 300 s; ALB allows in-flight
        requests to complete before removing a target.
\end{itemize}

\begin{sapdepth}
\textbf{SAP: Auto Scaling and ELB at enterprise scale:}
\begin{itemize}
  \item \textbf{Mixed Instances Policy}: combine On-Demand and Spot Instances in a single ASG;
        specify multiple instance types. AWS rebalances Spot allocations to minimise interruptions.
  \item \textbf{Warm pools}: maintain pre-initialised instances in a ``warm'' stopped state;
        reduce scale-out latency for apps with long boot times.
  \item \textbf{Scale-in protection}: mark specific instances as protected from scale-in
        (e.g.\ instances processing long jobs).
  \item \textbf{NLB + PrivateLink}: expose ALB or internal service to other AWS accounts/VPCs
        without VPC peering. NLB is the endpoint service; consumers create Interface Endpoints.
  \item \textbf{GWLB pattern}: VPC traffic $\rightarrow$ GWLB endpoint $\rightarrow$ GWLB $\rightarrow$
        3rd-party IDS/IPS (via GENEVE) $\rightarrow$ back to GWLB $\rightarrow$ destination.
        Deploy as a centralised inspection VPC (security VPC hub pattern).
\end{itemize}
\end{sapdepth}
